\section{Optimality of the kernel within the positivity method}\label{sec:optimality}

\subsection{Where the leading coefficient comes from}
In the proof of Theorem~\ref{thm:sharp} with a general admissible kernel $k=f*\tilde f$, $f\in C^2_c$ real,
$\operatorname{supp}f\subseteq[-a,a]$, only two terms are large. In the numerator the only term growing with
$q$ is $k(0)\log(q/\pi)$, coming from $\operatorname{Re}\Gamma_{q,\kappa}(k)$ through Plancherel. In the
denominator the only term growing with $x$ is $e^{y/2}\hat k(-\frac i2)=\sqrt x\int k(u)e^{u/2}du$. As $q\to\infty$
and $x/\log^2q\to\infty$ with $a$ fixed, the method therefore gives
\begin{equation}\label{eq:leading}
  \bias\le\bigl(\rho(f)+o(1)\bigr)\frac{\log q}{\sqrt x},\qquad
  \rho(f)=\frac{k(0)}{\int k(u)e^{u/2}\,du}=\frac{\norm f_2^2}{\langle f,e^{t/2}\rangle\,\langle f,e^{-t/2}\rangle}.
\end{equation}
The second expression follows from $k(0)=\norm f^2$ and $\hat k(z)=\hat f(z)\overline{\hat f(\bar z)}$ at $z=-\frac i2$.
For the bound to be meaningful we need both inner products of the same sign, which we may assume positive.

\subsection{The extremal problem}
\begin{lemma}[formal; \lean{GRHPseudo.inner_mul_inner_le}]\label{lem:ranktwo}
In a real inner product space, for all $f,\alpha,\beta$,
\[
  \langle f,\alpha\rangle\langle f,\beta\rangle\le\tfrac12\bigl(\norm\alpha\norm\beta+\langle\alpha,\beta\rangle\bigr)\norm f^2 .
\]
If $\norm\alpha\norm\beta>0$, equality holds iff $f$ is a multiple of $\alpha/\norm\alpha+\beta/\norm\beta$.
\end{lemma}
\begin{proof}
Put $u=\alpha/\norm\alpha$ and $v=\beta/\norm\beta$. Then
$4\langle f,u\rangle\langle f,v\rangle=\langle f,u+v\rangle^2-\langle f,u-v\rangle^2\le\norm f^2\norm{u+v}^2=2(1+\langle u,v\rangle)\norm f^2$.
Equality requires $f\perp u-v$ and $f\parallel u+v$; note that $u+v\perp u-v$.
\end{proof}

\begin{proof}[Proof of Proposition~\ref{prop:opt}]
Work in $L^2[-a,a]$ with $\alpha=e^{t/2}$ and $\beta=e^{-t/2}$. Then $\norm\alpha^2=\norm\beta^2=\int_{-a}^ae^t\,dt=2\sinh a$ and
$\langle\alpha,\beta\rangle=2a$. Lemma~\ref{lem:ranktwo} gives $\rho(f)\ge1/(\sinh a+a)$, with equality iff
$f\propto e^{t/2}+e^{-t/2}=2\cosh(t/2)$ on $[-a,a]$. The factorization statement is the theorem of Boas and Kac
and of Krein: a continuous positive-definite function supported in $[-2a,2a]$ is $f*\tilde f$ for some
$f\in L^2[-a,a]$ \cite{boaskac1945,ehmgneitingrichards2004}. Hence the bound covers every kernel to which
the positivity step applies, namely those with $\hat k\ge0$ on $\R$.
\end{proof}

\begin{proof}[Proof of Corollary~\ref{cor:asymp}]
\emph{Part 1.} A kernel supported in $[-2a,2a]$ and twisted by $y$ yields a measure supported in
$[xe^{-2a},xe^{2a}]$. With $X=xe^{2a}$ we have $\sqrt x=e^{-a}\sqrt X$, so \eqref{eq:leading} and
Proposition~\ref{prop:opt} give leading term $\ge c(a)\log q/\sqrt X$ with $c(a)=e^a/(a+\sinh a)$.

Minimizing: $(\log c)'(a)=1-(1+\cosh a)/(a+\sinh a)$ vanishes iff $a+\sinh a=1+\cosh a$, i.e.\
$a-1=\cosh a-\sinh a=e^{-a}$. There
\[
  c(a)=\frac{e^a}{1+\cosh a}=\frac{2e^a}{(e^{a/2}+e^{-a/2})^2}=\frac2{(1+e^{-a})^2}=\frac2{a^2}.
\]
Since $c(a)\to\infty$ as $a\to0^+$ and $c(a)\to2$ as $a\to\infty$, while $c(a_*)<2$, the critical point $a_*$ is the
global minimum.

\emph{Part 2.} Replace the cutoff $\eta_{a_*}$ by smooth cutoffs converging to $\mathbf 1_{[-a_*,a_*]}$ in $L^2$. This
makes $\rho(f)$ converge to $1/(a_*+\sinh a_*)$. For a fixed smooth kernel, every other term in
Theorem~\ref{thm:sharp} is $O_a(1)$ in the numerator, and in the denominator it is $O_a(1)+O(\log q\log x/\sqrt x)$,
which is $o(\sqrt x)$ once $x/\log^2q\to\infty$. The constant $c^{*2}=4/a_*^4=1.4972\ldots$ follows.
\end{proof}

\begin{remark}[Scope]
Proposition~\ref{prop:opt} concerns positive-definite kernels, those with $\hat k\ge0$. These are the kernels for which
the zero sum is controlled by positivity alone, with no zero counting; this is what makes our bounds explicit. The asymptotic
argument of \cite[Lemma~6.1]{lls2015} bounds the zero sum by $\sum_\gamma\abs{K(i\gamma)}$ instead, so it admits general kernels, with
leading coefficient $\frac1{2\pi}\int\abs{\hat k}/\int k(u)e^{u/2}du$. For such kernels we only have the trivial lower bound
$\ge1/(4\sinh a)$, from $\abs{k}\le\frac1{2\pi}\int\abs{\hat k}$. Whether a non-positive-definite kernel beats $c^*$ in that asymptotic
setting is open.
\end{remark}

\begin{remark}
Lemma~\ref{lem:ranktwo} is formalized; its specialization to $L^2[-a,a]$, the factorization theorem and the
asymptotic statement are not. The comparison with \cite{cmqr2024} is instructive: there the extremal problem is
one-sided, with sign conditions on $k$ adapted to a single non-residue condition, and the optimum is
$0.7615\ldots\log^2p$ for the least quadratic non-residue. Our problem bounds all characters simultaneously,
in both directions, using only $\hat k\ge0$.
\end{remark}
